\documentclass[10pt]{article}
\usepackage[margin=20mm]{geometry}
\usepackage{amsmath,amssymb,mathtools,fontspec,unicode-math,xurl,hyperref}
\IfFontExistsTF{DejaVu Serif}{\setmainfont{DejaVu Serif}}{\setmainfont{Latin Modern Roman}}
\setmathfont{latinmodern-math.otf}
\setlength{\parindent}{0pt}\setlength{\parskip}{6pt}\sloppy
\title{Finite-extension descent and sign checks for the cohomological Milnor formula}
\author{AI-assisted technical review}
\date{}
\begin{document}
\maketitle
Abstract: The acquired manuscript contains a false direct-sum assertion about finite Galois pushforward in the proof of Theorem 6.13. An explicit example over finite fields satisfies the stated hypotheses and distinguishes the two sheaves by Frobenius trace. The required characteristic-class scaling nevertheless follows from the projection formula and the manuscript’s own Lemma 3.2, with the analogous rank-twist property for characteristic cycles. The supported assessment is Major for the invalid finite-extension descent step in the proof of the main characteristic-class comparison. The explicit class-level replacement does not refute the theorem itself. Four admissible curve examples also reproduce the Milnor sign and its alternating-dimension convention. The higher-categorical and local Fourier foundations are identified as dependencies rather than independently certified proofs. Scientific adjudication remains pending.\par
\section*{Version, hypotheses, and inspected scope}
The source is Enlin Yang and Yigeng Zhao, Cohomological Milnor formula and Saito’s conjecture on characteristic classes, arXiv:2209.11086v4 [YZ]. Its arXiv margin is dated 16 February 2025 and the title-page manuscript date is 18 February 2025. The acquired PDF has 51 pages, with PDF and printed pagination agreeing. Only these author-manuscript bytes are assessed. The associated DOI is https://doi.org/10.1007/s00222-025-01319-y; equivalence with the publisher version is unverified.\par
Theorem 6.13 (p. 50), also stated as Theorem 1.3, proves the characteristic-class comparison for \(X\) smooth and quasi-projective over a perfect field \(k\) of characteristic \(p>0\). Coefficients form a finite local ring \(\Lambda\) whose residue characteristic differs from \(p\), and \(\mathcal F\) is constructible of finite Tor dimension. The broader ambient formulation of Conjecture 6.12 must not replace this proved scope.\par
Theorem 6.2 (p. 45) assumes a smooth curve \(Y\), a separated finite-type map \(f:X\to Y\), a closed point \(x\), and universal local acyclicity relative to \(\mathcal F\) on \(X\setminus\{x\}\). It states \[C^{\{x\}}_{X/Y/k}(\mathcal F)=-\operatorname{dimtot}R\Phi_{\bar x}(\mathcal F,f).\] It does not require \(X\) to be smooth. Theorem 6.5 gives the finite-exceptional-set fibration formula. Properness is required in the pushforward conductor formula, Theorem 6.7, but not in Theorem 6.5. The general non-isolated characteristic-cycle statement is explicitly Conjecture 6.10.\par
The helper manually read pp. 1–6, 14–21, 23–31, and 45–51, and actually viewed 17 selected original pages. The primary executor checked the decisive descent sentence, its internal repair, the non-acyclicity assumptions, and the Milnor/fibration statements against source images. Full text extraction is not claimed as full scientific verification. In particular, the complete transversality, specialization, and higher-categorical arguments were not independently reconstructed.\par
\section*{A local error in finite-extension descent}
Equation reference: eq:false. Problem 1 (Finite Galois pushforward need not be a trivial direct sum).  In the proof of Theorem 6.13, immediately before (6.13.2) on page 50, the manuscript asserts \[
 \sigma_*\sigma^*\mathcal F=\mathcal F^{\oplus d},
 \qquad d=[k':k],\] for the projection \(\sigma:X_{k'}\to X\) associated to a finite Galois extension of degree invertible in \(\Lambda\). This is false as an isomorphism of étale sheaves or derived objects.\par
The equality was confirmed in the original image; this is not an extraction ambiguity. Suppressed derived-functor symbols in the manuscript’s conventions do not resolve the issue, since the counterexample is already in degree zero and \(\sigma\) is finite étale.\par
\subsection*{A fully admissible example}
Take \[k=\mathbb F_5,\qquad k'=\mathbb F_{125},\qquad
 \Lambda=\mathbb F_2,\qquad X=\mathbb P^2_k,\qquad \mathcal F=\Lambda_X.\] The field \(k\) is perfect of positive characteristic; \(X\) is smooth, connected, projective, and therefore quasi-projective. The coefficient characteristic is distinct from five. The extension is finite Galois of degree three, which is invertible in \(\mathbb F_2\). The sheaf is constructible and flat, hence of finite Tor dimension. Every displayed condition of the descent paragraph is met.\par
Pull back [eq:false] (\#eq:false) to a \(k\)-rational point of \(X\). The stalk of \(\sigma_*\Lambda\) carries the permutation action on the three embeddings of \(k'\) into \(\bar k\). Arithmetic Frobenius acts as a three-cycle, represented after choosing an ordering by \[P=\begin{pmatrix}0&0&1\\1&0&0\\0&1&0\end{pmatrix}.\] The stalk of \(\Lambda^{\oplus3}\) instead has trivial Frobenius action \(I_3\). Over \(\mathbb F_2\), \[\operatorname{tr}(P)=0,\qquad \operatorname{tr}(I_3)=1.\] Thus the representations and sheaves cannot be isomorphic. The characteristic polynomials also differ: \[\det(tI-P)=t^3+1,\qquad
 \det(tI-I_3)=t^3+t^2+t+1.\] Trace is additive on exact sequences, so even interpreting [eq:false] (\#eq:false) as a Grothendieck-group equality of sheaves would not fix it.\par
An independent exact computation realizes the extension as \(\mathbb F_5[u]/(u^3+u+1)\). The cubic has values \(1,3,1,1,4\) on \(\mathbb F_5\), hence no root and is irreducible. Computation in this quotient gives \[u^5=1+u+4u^2,\quad
 u^{25}=4+3u+u^2,\quad u^{125}=u.\] It also verifies \(P^3=I\) and \(\operatorname{rank}(P-I)=2\) over \(\mathbb F_2\). These are actual finite-field and matrix calculations, not a purported computation of étale cohomology.\par
\subsection*{The required weaker assertion is available}
Equation reference: eq:correct. The projection formula gives the correct identity \[
 \sigma_*\sigma^*\mathcal F\simeq\mathcal F\otimes\mathcal L,
 \qquad \mathcal L=\sigma_*\Lambda,\] where \(\mathcal L\) is locally constant, constructible, and flat of rank \(d\). It is a permutation local system and need not be constant.\par
Equation reference: eq:rank. The manuscript’s Lemma 3.2, equation (3.2.1) on page 15, proves exactly the required cohomological rank-twist identity: \[
 C_{X/k}(\mathcal F\otimes\mathcal L)=d\,C_{X/k}(\mathcal F).\] For disconnected \(X\), apply it componentwise. Therefore the last degree-scaling equality in (6.13.2) remains justified without [eq:false] (\#eq:false).\par
On the characteristic-cycle side, étale locality and additivity give \[\operatorname{CC}(\mathcal F\otimes\mathcal L)=d\,\operatorname{CC}(\mathcal F).\] Indeed, after a surjective étale cover trivializing \(\mathcal L\), the assertion is direct-sum additivity. Equality of cycle coefficients descends since pullback along a surjective étale cover detects them. Applying the zero-section construction of \(\operatorname{cc}\) gives \(\operatorname{cc}_X(\mathcal F\otimes\mathcal L)=d\,\operatorname{cc}_X(\mathcal F)\). This uses the usual characteristic-cycle foundations already underlying the paper; it does not assume the desired cohomological comparison and is not circular.\par
Thus both outer degree-scaling equalities in (6.13.2) have a direct replacement. The middle steps still use the stated comparison over \(k'\), proper pushforward, and the cited characteristic-cycle input. Since \(d\) is invertible in \(\Lambda\), the intended descent conclusion follows once those stated dependencies are granted.\par
\subsection*{Impact and repair}
The supported assessment is Major for the printed core derivation. The false sheaf identity is invoked within the stated hypotheses of Theorem 6.13 to justify both outer degree-scaling equalities in (6.13.2); it is not an unused assertion or a boundary excluded from the theorem. Its short replacement at the characteristic-class level does not validate that sheaf identity. The demonstrated defect concerns this descent step, and there is no counterexample here to the conclusion of Theorem 6.13 or to the Milnor formula. The repair is to replace [eq:false] (\#eq:false) by [eq:correct] (\#eq:correct), identify the rank-\(d\) permutation local system, and cite Lemma 3.2 together with the rank-twist property of characteristic cycles. If the original sentence intended only equality after applying characteristic classes, that weaker meaning should be stated explicitly.\par
This report does not assert that every other step of the general proof has been certified. Nor does it assert that a later publisher version contains the same sentence or that the correction is previously unknown.\par
\section*{Definitions and proof dependencies checked}
Section 4.1 (p. 23) distinguishes the following assumptions. Condition (C1) controls the locally constant relative dualizing object for \(Y/S\). Condition (C3) imposes finite Tor dimension and universal local acyclicity away from \(Z\) over \(Y\) and on \(X\) over \(S\). These define the tilded non-acyclicity class. The untilded class additionally requires \[\text{(C2)}\qquad H^0(Z,\mathcal K_{Z/Y})=H^1(Z,\mathcal K_{Z/Y})=0.\] The triangle on page 24 induces \[H^0_Z(\mathcal K_{X/Y})\longrightarrow
 H^0_Z(\mathcal K_{X/S})\longrightarrow
 H^0_Z(\mathcal K_{X/Y/S})\longrightarrow
 H^1_Z(\mathcal K_{X/Y}).\] Under (C2), the outer groups vanish and the middle arrow is an isomorphism. This verifies the unique lifting mechanism with its actual hypotheses. For a closed point over a smooth curve, purity gives \(\mathcal K_{x/Y}=\Lambda(-1)[-2]\), up to the finite étale residue-field step, so the requisite \(H^0,H^1\) vanish. This concerns \(x\to Y\) and does not impose smoothness on \(X\).\par
The proof of Theorem 6.2 uses étale locality, the Artin–Schreier Fourier deformation, a universally locally acyclic extension, specialization, smooth base change, a projection formula with a flat local system, and the cited \(\Psi\)-goodness result (pp. 45–47). The final point-supported trace is an alternating rank. The elementary rank calculation does not replace these foundations.\par
Theorem 6.5 combines the fibration identity of Theorem 5.1 with isolated point terms. Theorem 6.7 uses proper pushforward and correctly retains the tilded class on a possibly positive-dimensional exceptional fiber. For Theorem 6.13, good fibrations and rational differentials reduce the comparison to lower-dimensional fibers and isolated vanishing cycles; the blowup identities in Lemma 3.15 provide the induction step after the cited good-pencil construction. These uses and their scopes were traced, but the full cited proofs and higher-categorical construction were not independently rebuilt.\par
\section*{Four sign and boundary checks}
For these checks take \(k=\overline{\mathbb F}_5\) and \(\Lambda=\mathbb F_{11}\). The coefficient characteristic is different from five, and \(\Lambda^\times\) contains a nontrivial fifth root of unity, so the additive-character choice in the Fourier argument poses no issue for these examples. All sheaves below are constructible of finite Tor dimension.\par
For \(X=Y=\mathbb P^1_k\), \(f=\mathrm{id}\), \(x=0\), and immersions \(i:\{x\}\hookrightarrow X\), \(j:X\setminus\{x\}\hookrightarrow X\), the vanishing-cycle triangle is \[i^*\mathcal F\longrightarrow R\Psi\mathcal F\longrightarrow R\Phi\mathcal F
 \longrightarrow i^*\mathcal F[1].\] The total dimension is alternating in cohomological degree. A copy of \(\Lambda[1]\) contributes \(-1\), not \(+1\).\par
\subsection*{The identity map}
For \(\mathcal F=\Lambda_X\), specialization is the identity, so \(R\Phi=0\) and the non-acyclicity contribution is zero.\par
For \(\mathcal F=i_*\Lambda\), the generic restriction and nearby cycles vanish. Thus \(R\Phi=\Lambda[1]\), \(\operatorname{dimtot}R\Phi=-1\), and the proposed point contribution is \(+1\). Independently, proper pushforward of the point gives \(C_{X/k}(i_*\Lambda)=[x]\), while its relative class for the identity has generic rank zero. The already established fibration identity therefore has correction \(+[x]\). On \(\mathbb P^1\), this point class is a nonzero degree-one generator, so the coefficient is detected.\par
For \(\mathcal F=j_!\Lambda\), the special stalk is zero and nearby cycles are rank one in degree zero with trivial inertia. Hence \(\operatorname{dimtot}R\Phi=1\) and the predicted correction is \(-1\). The triangle \(j_!\Lambda\to\Lambda_X\to i_*\Lambda\) and additivity give \[C_{X/k}(j_!\Lambda)=c_1(T_X)-[x].\] The relative rank is one, so the correction from the fibration identity is exactly \(-[x]\). The signs \(+1\) and \(-1\) are distinct in \(\mathbb F_{11}\).\par
\subsection*{A tame map with two branch points}
Take \(f:\mathbb P^1\to\mathbb P^1\), \([u:v]\mapsto[u^2:v^2]\), and \(\mathcal F=\Lambda_X\). It is finite and étale outside \(0,\infty\); the ramification index two is prime to five. At each branch point the nearby fiber has two points and specialization is the diagonal map \(\Lambda\to\Lambda^2\). The cokernel has rank one and Swan conductor zero. Thus each vanishing-cycle total dimension is \(1\).\par
The absolute class has degree \(\deg c_1(T_{\mathbb P^1})=2\), while the pulled-back target tangent class has degree \(4\). The fibration identity therefore needs total correction \(2-4=-2\), agreeing with two local contributions \(-1\). Coordinate inversion exchanges the branch points and preserves the map, so their local coefficients agree; two is invertible in \(\mathbb F_{11}\). The numerical degree balance is \[4-1-1=2.\] The accompanying script verifies the Euler-rank arithmetic and finite matrices. It does not implement vanishing cycles or the six operations. The geometric arguments above supply the content beyond the arithmetic.\par
These checks support the stated minus sign and the alternating total dimension. They do not verify arbitrary wild ramification, all constructible complexes, or the conjectural non-isolated case.\par
\section*{Disposition}
The invalid descent step in the proof of Theorem 6.13 warrants the Major assessment under the core-derivation criterion. It should be corrected using the existing rank-twist lemma and characteristic-cycle locality. No failure of the main characteristic-class equality is established. The principal restrictions, lifting vanishing conditions, and tested Milnor signs are consistent.\par
The complete transversality and higher-categorical constructions, the local Fourier and universally locally acyclic extension theorems, coefficient-enlargement compatibility in the general Fourier argument, and the external characteristic-cycle and good-pencil inputs remain outside the independent proof scope. They are not promoted to confirmed errors. No author code was executed, and absence of code is not treated as a mathematical defect. Scientific decision is pending, `human\_reviews` is empty, and the independent focused check is recorded separately from that pending scientific adjudication.\par
\section*{Supporting checks}
The severity and impact statement have been revised; the current build and actual final-page inspection are bound to the revised TeX and PDF in provenance. Runtime self-description and controller UI observations remain separate, with backend identity unverified.\par
The first helper computation failed during JSON serialization and is preserved with its failure record. A distinct corrected script actually completed with exit code zero; only that result is counted as successful. Source images, separate render/view receipts, complete hypotheses and derivations, source PDF/manifest, scripts, build sources/logs, and actual support-file hashes accompany this report.\par
1 E. Yang and Y. Zhao, Cohomological Milnor formula and Saito’s conjecture on characteristic classes, arXiv:2209.11086v4, 16 February 2025; manuscript dated 18 February 2025, 51 pp. https://arxiv.org/abs/2209.11086v4.\par
\end{document}
